% Macros used in the paper					% -*- latex -*-

\makeatletter

\newcommand{\ie}{\mbox{\emph{i.e.}}}
\newcommand{\Ie}{\mbox{\emph{I.e.}}}
\newcommand{\eg}{\mbox{\emph{e.g.}}}
\newcommand{\Eg}{\mbox{\emph{E.g.}}}
\newcommand{\etc}{\emph{etc.}}
\newcommand{\nb}{\mbox{\emph{n.b.}}}
\newcommand{\Nb}{\mbox{\emph{N.b.}}}

\newcommand{\sequence}[1]{#1 \, \ldots \,}
\newcommand{\llb}{\left\llbracket}
\newcommand{\rrb}{\right\rrbracket}

\newcommand{\mdoubleplus}{\mathbin{+\mkern-7mu+}}
\newcommand{\thappend}{\mdoubleplus}
\newcommand{\thempty}{\square}
\newcommand{\thprefix}{\sqsubseteq}
\newcommand{\theq}{\equiv}
\providecommand{\dotdiv}{% Don't redefine it if available
  \mathbin{% We want a binary operation
    \vphantom{+}% The same height as a plus or minus
    \text{% Change size in sub/superscripts
      \mathsurround=0pt % To be on the safe side
      \ooalign{% Superimpose the two symbols
        \noalign{\kern-.35ex}% but the dot is raised a bit
        \hidewidth$\smash{\cdot}$\hidewidth\cr % Dot
        \noalign{\kern.35ex}% Backup for vertical alignment
        $-$\cr % Minus
      }}}}
\newcommand{\thmonus}{\dotdiv}

% Syntactic classes
%\newcommand{\scal}{\mathit{s}}
\newcommand{\term}{\mathit{t}}
\newcommand{\atom}{\mathfrak{a}}
\newcommand{\func}{\mathfrak{f}}
\newcommand{\atval}{\mathfrak{v}}
\newcommand{\expr}{\mathit{e}}
\newcommand{\var}{\mathit{x}}
\newcommand{\primop}{\mathfrak{o}}
\newcommand{\val}{\mathit{v}}
\newcommand{\baseval}{\mathfrak{b}}
\newcommand{\type}{\mathit{\tau}}
\newcommand{\basetype}{\mathit{B}}
\newcommand{\kind}{\mathit{k}}
\newcommand{\idx}{\mathit{\iota}}
\newcommand{\sort}{\mathit{\gamma}}
\newcommand{\nat}{\mathit{n}}
\newcommand{\tEnv}{\Gamma}
\newcommand{\kEnv}{\Delta}
\newcommand{\sEnv}{\Theta}
\newcommand{\emptyEnv}{\cdot}
\newcommand{\ctxt}{\mathbb{E}}
\newcommand{\hole}{\square}
\newcommand{\erased}[1]{\widehat{#1}}
\newcommand{\eexpr}{\erased{\expr}}
\newcommand{\eatom}{\erased{\atom}}
\newcommand{\eterm}{\erased{\term}}
\newcommand{\efunc}{\erased{\func}}
\newcommand{\evalue}{\erased{\val}}
\newcommand{\eatvalue}{\erased{\atval}}
\newcommand{\ectxt}{\erased{\ctxt}}

\newcommand{\Expr}{\mathit{Expr}}
\newcommand{\Atom}{\mathit{Atom}}
\newcommand{\Func}{\mathit{Func}}
\newcommand{\Primop}{\mathit{Op}}
\newcommand{\Term}{\mathit{Term}}
\newcommand{\Type}{\mathit{Type}}
\newcommand{\Idx}{\mathit{Idx}}
\newcommand{\Kind}{\mathit{Kind}}
\newcommand{\Sort}{\mathit{Sort}}
\newcommand{\Val}{\mathit{Val}}
\newcommand{\Atval}{\mathit{Atval}}
\newcommand{\Ctxt}{\mathit{Ctxt}}
\newcommand{\wt}[1]{{#1}^{T}}

% Signature, specifying primops' types
\newcommand{\Signature}{\mathcal{S}}
\newcommand{\Sigref}[1]{\Signature{\hskip-.25em}\llb#1\rrb}

% Evaluation context filled with term
\newcommand{\fillc}[2]{#1\hskip-.15em\left[#2\right]}

% Tools for sexp-like things
\newcommand{\parens}[1]{\left( #1 \right)}
\newcommand{\ttparens}[1]{\text{\tt (}#1\text{\tt )}}
\newcommand{\braces}[1]{\left\{ #1 \right\}}
\newcommand{\ttbraces}[1]{\text{\tt \{}#1\text{\tt \}}}
\newcommand{\bracks}[1]{\left[ #1 \right]}
\newcommand{\ttbracks}[1]{\text{\tt [}#1\text{\tt ]}}
\newcommand{\pseq}[1]{\parens{\sequence{#1}}}
% For sexp forms, give the keyword as the first argument and the rest
% of the subexprs as the second argument
\newcommand{\sexp}[2]{\left( #1 \subexp{#2} \right)}
\newcommand{\ttsexp}[2]{\text{\tt(}#1 \subexp{#2}\text{\tt)}}
\newcommand{\ttseq}[1]{\ttsexp{#1}{\text{\tt ...}}}
%\newcommand{\noarg}{\negthickspace}
% If a subexp might be blank, don't necessarily add spacing to separate it
\newcommand{\subexp}[1]{\ifblank{#1}{}{\;#1}}
% Variant for possibly-missing subexp at beginning of list
\newcommand{\subexpcar}[1]{\ifblank{#1}{}{#1\;}}

\let\oldl\l %Save the old \l on \oldl
\newcommand{\cl}{\llap{$\lambda$\hskip-.65em}\phantom{ss}}
%% JFP-friendly versions of \lambda-in-code macros
\renewcommand{\l}[1]{\cl(#1)}
\newcommand{\ttlm}[1]{\cl\ttparens{#1}}
%% Non-JFP versions of \lambda-in-code macros
% \renewcommand{\l}[1]{\ \llap{$\lambda$\hskip-.5em}\ (#1)}
% \newcommand{\ttlm}[1]{\ \llap{$\lambda$\hskip-.5em}\ \ttparens{#1}}

% Syntactic forms
\newcommand{\annotate}[2][]{{#2}^{#1}}
\newcommand{\notevar}[2]{\ttparens{#1 \; #2}}

%\newcommand{\tuple}[2]{\langle #1 , #2 \rangle}
\newcommand{\tuple}[1]{\ttsexp{\text{\tt tuple}}{#1}}
\newcommand{\notedtuple}[2]{\annotate[#1]{\tuple{#2}}}

%\newcommand{\lam}[2]{\parens{\lambda \; \bracks{#1} \; #2}}
\newcommand{\lam}[2]{\ttsexp{\ttlm{#1}}{#2}}
\newcommand{\notedlam}[3][]{\annotate[#1]{\lam{#2}{#3}}}

% \newcommand{\tlam}[2]{
%   \parens{{\tt T\lambda} \bracks{#1} \; #2}}
\newcommand{\tlam}[2]{\ttsexp{{\tt T}\ttlm{#1}}{#2}}
\newcommand{\notedtlam}[3][]{\annotate[#1]{\tlam{#2}{#3}}}

% \newcommand{\ilam}[2]{
%   \parens{{\tt I\lambda} \bracks{#1} \; #2}}
\newcommand{\ilam}[2]{\ttsexp{{\tt I}\ttlm{#1}}{#2}}
\newcommand{\notedilam}[3][]{\annotate[#1]{\ilam{#2}{#3}}}

%\newcommand{\frm}[2]{\text{\tt Frm}\bracks{#1}_{#2}}
%\newcommand{\frm}[2]{\ttsexp{\text{\tt frame}}{\ttparens{#2} \ifblank{#1}{}{\;#1}}}
\newcommand{\frm}[2]{\ttsexp{\text{\tt frame}}{\ttparens{#2} \subexp{#1}}}
\newcommand{\notedfrm}[3][]{\annotate[#1]{\frm{#2}{#3}}}

%\newcommand{\emptyfrm}[2]{\text{\tt Frm}^{#1}_{#2}}
\newcommand{\emptyfrm}[2]{\ttsexp{\text{\tt frame}}{\ttparens{#2} \subexp{#1}}}
\newcommand{\notedemptyfrm}[3][]{\annotate[#1]{\frm{#2}{#3}}}

%\newcommand{\arrlit}[2]{\text{\tt Arr}\bracks{#1}_{#2}}
\newcommand{\arrlit}[2]{\ttsexp{\text{\tt array}}{\ttparens{#2} \subexp{#1}}}
\newcommand{\notedarrlit}[3][]{\annotate[#1]{\arrlit{#2}{#3}}}

\newcommand{\emptyarrlit}[2]{\ttsexp{\text{\tt array}}{\ttparens{#2} \subexp{#1}}}
\newcommand{\notedemptyarrlit}[3][]{\annotate[#1]{\frm{#2}{#3}}}

%\newcommand{\app}[2]{\parens{#1 \; #2}}
\newcommand{\app}[2]{\ttsexp{#1}{#2}}
\newcommand{\notedapp}[3][]{\annotate[#1]{\app{#2}{#3}}}

%\newcommand{\tapp}[2]{\parens{\text{\tt T-APP} \; #1 \; #2}}
\newcommand{\tapp}[2]{\ttsexp{\text{\tt t-app}}{#1 \subexp{#2}}}
\newcommand{\notedtapp}[3][]{\annotate[#1]{\tapp{#2}{#3}}}

\newcommand{\iapp}[2]{\ttsexp{\text{\tt i-app}}{#1 \subexp{#2}}}
\newcommand{\notediapp}[3][]{\annotate[#1]{\iapp{#2}{#3}}}

%\newcommand{\dsum}[3]{\parens{ \text{\tt box} \; #1 \; #2 \; : #3 }}
\newcommand{\dsum}[3]{\ttsexp{\text{\tt box}}{#1 \subexp{#2} \subexp{#3}}}
\newcommand{\noteddsum}[3][]{
  \annotate[#1]{
    \ttsexp{ \text{\tt box} \subexp{#2} \subexp{#3}}}}

\newcommand{\dproj}[4]{
  \ttsexp
  {\text{\tt unbox}}
  % Typesetting this is weird because the "optional" piece is #1, so its
  % trailing space should be guarded by \ifblank{#1}, rather than considering it
  % attached to #2
  {\ttparens{\subexpcar{#1}#2 \subexp{#3}} \subexp{#4}}}
\newcommand{\noteddproj}[5][]{\annotate[#1]{\dproj{#2}{#3}{#4}{#5}}}

\newcommand{\typefun}[2]{\ttsexp{\text{\tt ->}}{\ttparens{#1} \; #2}}
%{\ttsexp{#1}{\text{\tt ->} \; #2}}
\newcommand{\typeprod}[2]{\ttsexp{\text{\tt Tuple}}{#1 \; #2}}
\newcommand{\typeuniv}[2]{\ttsexp{\text{\tt Forall}}{\ttparens{#1} \; #2}}
\newcommand{\typedprod}[2]{\ttsexp{\text{\tt Pi}}{\ttparens{#1} \; #2}}
\newcommand{\typedsum}[2]{\ttsexp{\text{\tt Sigma}}{\ttparens{#1} \; #2}}
%\newcommand{\typearray}[2]{#1 #2}
\newcommand{\typearray}[2]{\ttsexp{\text{\tt Arr}}{#1 \; #2}}

\newcommand{\idxshape}[1]{\ttsexp{\text{\tt Shp}}{#1}}
\newcommand{\idxadd}[1]{\ttsexp{\text{\tt +}}{#1}}
\newcommand{\idxappend}[1]{\ttsexp{\text{\tt ++}}{#1}}
%\newcommand{\idxappend}[1]{\braces{\thappend \; #1}}

\newcommand{\kindarray}{\text{\tt Array}}
\newcommand{\kindatom}{\text{\tt Atom}}

\newcommand{\sortdim}{\text{\tt Dim}}
\newcommand{\sortshp}{\text{\tt Shape}}


% Judgment forms
\newcommand{\rulename}[1]{\quad \textsc{#1}}
\newcommand{\seqpremise}[2]{#2 \; \; \text{for each $#1$}}

\newcommand{\hassort}[2]{#1 :: #2}
\newcommand{\sortof}[3]{#1 \vdash \hassort{#2}{#3}}

\newcommand{\haskind}[2]{#1 :: #2}
\newcommand{\kindof}[4]{#1;#2 \vdash \haskind{#3}{#4}}

\newcommand{\hastype}[2]{#1 : #2}
\newcommand{\typeof}[5]{#1;#2;#3 \vdash \hastype{#4}{#5}}

\newcommand{\typeeqsym}{\cong}
\newcommand{\teqv}[2]{#1 \typeeqsym #2}
\newcommand{\ieqv}[2]{\text{\sc Valid} \llb #1 \theq #2 \rrb}

\newcommand{\wfenv}[3]{#1 ; #2 \vdash #3}

\newcommand{\fresh}[1]{\quad\text{with fresh $#1$}}


% "Replace Y with Z"
\newcommand{\singlesubst}[2]{#1 \mapsto #2}
% "In X, replace ____"
\newcommand{\multisubst}[2]{#1\hspace{-2pt}\bracks{#2}}
% "In X, replace Y with Z
\newcommand{\subst}[3]{\multisubst{#1}{\singlesubst{#2}{#3}}}
% "In X, replace Ys with Zs"
\newcommand{\seqsubst}[3]{\multisubst{#1}{\sequence{\singlesubst{#2}{#3},}}}


% Evaluation rules
\newcommand{\evalrule}[4][]{
  \mbox{
    \begin{tabular}{l}
       #2 \\ $ \; \mapsto_{\mathit{#1}} $ #3
    \end{tabular}
  }\\
  \mbox{#4}
}
\newcommand{\rulesep}{\medskip\\}
\newcommand{\0}{\hspace*{0em} }


% Proof utilities
\newcommand{\pfstub}{{\color{red} Left as an exercise for the author}}
\newcommand{\cfitem}[4][]{
\item{
  If $\type$ is of the form $#3$,\\
  then $#2$ is of the form $#4$#1.}
}
\newcommand{\cfarritem}[4][]{
  \cfitem
      [#1]
      {#2}
      {\typearray{#3}{\idx}}
      {\arrlit{\sequence{#4}}{\sequence{\nat}}}}

% "Collapsible" proof environment: give \pfskip as argument to elide the proof
\newcommand{\pfskip}[1]{{\color{cyan} elided for editing}}
\NewEnviron{cproof}[1][]{\begin{proof}#1{\BODY}\end{proof}}
% "Sketchable" proof environment: consults \pftype to pick version of proof
% (can select \fullproof or \sketchproof or \collapseproof)
\def\fullproof{0}
\def\sketchproof{1}
\def\collapseproof{2}
\NewEnviron{sproof}[1][{\color{magenta} sketch missing}]
{\begin{proof}
\if\pftype\fullproof {\BODY} \else\if\pftype\sketchproof #1 \else \pfskip{} \fi\fi
\end{proof}}
\newcommand{\TODO}{{\color{red} [TODO]}}

% Gadgets for sketching judgment derivations
\newcommand{\somehow}[1]{\inferrule*{\vdots}{#1}}
\newcommand{\teqvInd}[2]{\somehow{\teqv{#1}{#2}}}
\newcommand{\typeofInd}[5]{\somehow{\typeof{#1}{#2}{#3}{#4}{#5}}}
\newcommand{\kindofInd}[4]{\somehow{\kindof{#1}{#2}{#3}{#4}}}
\newcommand{\sortofInd}[3]{\somehow{\sortof{#1}{#2}{#3}}}
\newcommand{\infr}[3][]{\inferrule*[right=\;#1]{#2}{#3}}
\newcommand{\ctrand}{\[\text{and}\]}

% Erasure functions and related things
\newcommand{\EraseE}[1]{\mathcal{E}{\hskip-.25em}\llb #1 \rrb}
\newcommand{\EraseA}[1]{\mathcal{A}{\hskip-.25em}\llb #1 \rrb}
\newcommand{\EraseT}[1]{\mathcal{T}{\hskip-.25em}\llb #1 \rrb}
\newcommand{\EraseC}[1]{\mathcal{C}{\hskip-.25em}\llb #1 \rrb}
\newcommand{\EraseTrm}[1]{\mathcal{R}{\hskip-.25em}\llb #1 \rrb}
\newcommand{\EraseEqv}{\cong_{\mathcal{E}}}

%% What you frequently want when you say \tt:
\def\ttt{\ttfamily\upshape\mdseries\catcode``=13\@noligs\frenchspacing}

%% Centered version of alltt environment
\newenvironment{centertt}{%
  \par\medskip
  \centering
  \varwidth{\linewidth}%
  \alltt
}{%
  \endalltt
  \endvarwidth
  \vspace{-5pt}
  \par
}

% Works in math mode; all special chars remain special; cheaper than \cd.
% Will not be correct size in super and subscripts, though.
\newcommand{\ex}[1]{\mbox{\ttt{#1}}}

\newenvironment{inset}
 {\bgroup\parskip=1ex plus 1ex\begin{list}{}%
        {\topsep=0pt\rightmargin\leftmargin}%
        \item[]}%
 {\end{list}\leavevmode\egroup\global\@ignoretrue}

\newenvironment{tightinset}
 {\bgroup\parskip=0pt\begin{list}{}%
        {\topsep=0pt\rightmargin\leftmargin}%
        \item[]}%
 {\end{list}\leavevmode\egroup\ignorespaces}

\newenvironment{leftinset}
 {\bgroup\parskip=1ex plus 1ex\begin{list}{}%
        {\topsep=0pt}%
        \item[]}%
 {\end{list}\leavevmode\egroup\global\@ignoretrue}

% For typesetting interactive examples.
\newcommand{\result}[1]{\textsl{#1}}

% For multiletter vars in math mode:
\renewcommand{\v}[1]{\ensuremath{\mathit{#1}}} % \frenchspacing?
\newcommand{\vi}[2]{\ensuremath{\mathit{#1}_{#2}}}

\newcommand{\tu}{\ensuremath{\rightarrow}}

\newcommand{\ft}[2]{(\tu\ (#1) #2)}
\newcommand{\faty}[2]{(\(\forall\) (#1) {#2})}
\newcommand{\pity}[2]{(\(\Pi\) (#1) {#2})}
\newcommand{\vari}[2]{\ensuremath{#1_{#2}}}

\makeatother
